\section{Detection and Attribution on Real LLM Deliverables}
\label{sec:robustness}

\subsection{Unedited leaks: the operating point}
Table~\ref{tab:ea0} gives the no-attack round trip on 100 held-out Chinese
deliverables with a single-user registry (uid 7): existence and attribution
are both 100\%, calibrated and uncalibrated. Real LLM deliverables sit
comfortably inside the watermark's capability window---their function-word
and connector vocabulary (``not only\ldots but also'', ``therefore'',
``should'') is precisely what zero-cost dictionaries cover.

\begin{table}[t]
\centering\small
\caption{E-A0: unedited-leak round trip on 100 real LLM Chinese
deliverables (zero-cost codec, explicit metadata path).}
\label{tab:ea0}
\begin{tabular}{@{}lccc@{}}
\toprule
Configuration & Existence TPR & Attribution exact & $n$ \\
\midrule
Uncalibrated defaults & 1.00 & 1.00 & 100 \\
In-domain calibrated & 1.00 & 1.00 & 100 \\
\bottomrule
\end{tabular}
\end{table}

\subsection{Real LLM paraphrase: the capability boundary}
We attack the marked texts with the same 32B LLM at three rewrite strengths
(``replace $\sim$15\%/30\%/50\% of content words, keep meaning and
length''; $n{=}50$ per strength, calibrated and uncalibrated).
Table~\ref{tab:ea2} shows the outcome: at 15\% rewrite strength, existence
TPR is 2\% and attribution is zero; at 30\% and 50\%, both are zero
everywhere. Real LLM paraphrase is far more destructive than the synthetic
within-group synonym attacks used in prior robustness reports, because the
rewriter varies exactly the connector and function words that carry the
mark, not just content words.

\begin{table}[t]
\centering\small
\caption{E-A2: survival under real LLM paraphrase (Qwen2.5-32B rewriter,
$n{=}50$ per cell; zero-cost Chinese codec). This is a capability boundary,
not a robustness table: semantic rewriting defeats lexical watermarks.}
\label{tab:ea2}
\begin{tabular}{@{}lccc@{}}
\toprule
Rewrite strength & Existence TPR & Attribution & Abstain \\
\midrule
15\% & 0.02 & 0.00 & 0.00 \\
30\% & 0.00 & 0.00 & 0.00 \\
50\% & 0.00 & 0.00 & 0.00 \\
\bottomrule
\end{tabular}
\end{table}

We therefore state the capability window plainly: AAWM traces
\emph{unedited or structurally edited} leaks (insertions, deletions,
reorderings, excerpting survive by design through content-addressed
voting). Against an attacker who semantically rewrites the leaked document
with an LLM, no lexical watermark survives; detection of the mark's
remains is a research problem, and the abstain protocol's honest output is
``watermark destroyed''. Deployments that expect rewrite attacks should
combine watermarking with traditional forensics rather than rely on any
lexical scheme.
